\documentclass[12pt,a4paper]{article}
\usepackage[a4paper,twoside,top=2.5cm,bottom=2cm,inner=3.5cm,outer=2cm]{geometry}
\usepackage{setspace}
\usepackage[hidelinks]{hyperref}

\usepackage{fontspec} % Dla Arial na stronie tytułowej
\newfontfamily\titlepagefont{Arial} % Ustalenie czcionki dla strony tytułowej
%\setstretch{1.5} 


\usepackage{newtxtext} % Czcionka New Times Roman
\usepackage{graphicx}
\usepackage{adjustbox}
\usepackage[utf8]{inputenc}
\usepackage[english, polish]{babel}
\babelprovide[transforms = oneletter.nobreak]{polish} 
\usepackage[T1]{polski}
\usepackage{array}
\usepackage{enumitem}
%\usepackage{biblatex}
\usepackage{tabularx}
\usepackage{xltabular}
\usepackage{multirow}
\usepackage{stackengine}
\usepackage{tocloft}
\usepackage{enumitem}
\usepackage{chngcntr}
\usepackage{newfloat}
\usepackage[font=normalsize]{caption}
\usepackage[justification=centering]{caption}
\usepackage{boldline}
\usepackage{indentfirst}
\usepackage{listings}
\usepackage{csquotes}
\usepackage{float}
\usepackage{capt-of}
\usepackage[backend=biber, style=numeric]{biblatex}



\DeclareFloatingEnvironment{graph}

\addbibresource{bib/thesis.bib}

\newcolumntype{C}[1]{>{\centering\arraybackslash} m{#1} }
\newcolumntype{?}{!{\vrule width 1.5pt}}

\addto\captionsenglish{
  \renewcommand{\figurename}{Rysunek}
  \renewcommand{\thefigure}{\arabic{section}.\arabic{figure}}
  \renewcommand{\graphname}{Listing}
  \renewcommand{\thegraph}{\arabic{section}.\arabic{graph}}
  \renewcommand{\contentsname}{Spis treści}
}


\addto\captionsenglish{
    \renewcommand{\figurename}{Rysunek}
    \renewcommand{\graphname}{Listing}
    \renewcommand{\thegraph}{\arabic{section}.\arabic{graph}}
    \renewcommand{\contentsname}{Spis treści}
}


\renewcommand{\thefigure}{\arabic{section}.\arabic{figure}}

\newcommand{\row}[1]{\hline #1 \\}
\newcommand{\image}[2]{\begin{figure}[h!]
    \centering
    \includegraphics[width=\linewidth]{#1}
    \caption{#2}
    \label{fig:image}
\end{figure}
\par}

\newcommand{\lis}[3]{\begin{graph}[h!]
    \centering
    \caption{#2}
    \includegraphics[#3]{#1}
    \label{fig:graph}
\end{graph}
\par}

\def\blankpage{%
      \clearpage%
      \thispagestyle{empty}%
      \null%
      \clearpage}

\lstset{
    basicstyle=\small\ttfamily, % Adjust the font size here
   % framerule=0.5pt,
}

\floatstyle{plain}
\newfloat{listing}{htbp}{lop}[section]
\floatname{listing}{Listing}

\captionsetup[listing]{position=top}

\newcommand{\newCharapter}[1]{\newpage \include{#1}}


\begin{document}

\counterwithin{figure}{section}
\counterwithin{lstlisting}{section}
\counterwithin{table}{section}

\begin{titlepage}

    \titlepagefont
    \hspace*{-0.5cm}{\includegraphics[width=7cm]{img/title/weii-2.jpg}} \\
    \vspace{4cm}
    \break
    \fontsize{1.4cm}{1.4cm}\titlepagefont {Praca dyplomowa} \\
    \fontsize{1.4cm}{1.4cm}\titlepagefont {inżynierska} \\
    \vspace{0.5cm}
    \break
    \normalsize{} na kierunku Informatyka \\
    w specjalności Inżynieria Oprogramowania \\
    \vspace{1cm}
    \break
    \large
    Aplikacja mobilna wspomagająca turystykę górską. \\
    \break
    Mobile application supporting mountain tourism. \\
    \vspace{1.5cm}
    \break
    Magdalena Kamińska \\
    97628   \\
    \vspace{0.5cm}
    \break
    Katarzyna Jakubowska \\
    97620   \\
    \vspace{1.5cm}
    \break
    \normalsize{}
    Promotor Dr inż. Piotr Kopniak
    \vspace*{\fill}
    \vspace{0.5cm}
    \begin{figure*}[bbp]
        {\titlepagefont Lublin 2025}
    \end{figure*}
\end{titlepage}
\newpage

\onehalfspacing
\blankpage

\tableofcontents
\newCharapter{chapters/chapter0}
\setcounter{figure}{0}
\newCharapter{chapters/chapter1}
\setcounter{figure}{0}
\newCharapter{chapters/chapter2}
\setcounter{figure}{0}
\newCharapter{chapters/chapter3}
\setcounter{figure}{0}
\newCharapter{chapters/chapter5}
\setcounter{figure}{0}
%\newCharapter{chapters/chapter4}
%\setcounter{figure}{0}
\newCharapter{chapters/chapter6}
\setcounter{figure}{0}
\newCharapter{chapters/testy}
\setcounter{figure}{0}
\newCharapter{chapters/chapter7}

%\listoffigures

%\printbibliography[type=book, title={Źródła naukowe}]
%\printbibliography[type=online,title={Źródła internetowe}]
\printbibliography[sorting=none]


\end{document}
